\chapter{Michael Francis Atiyah --- Fields Medal 1966}\label{ch1966michaelfrancisatiyah:michael-francis-atiyah-fields-medal-1966-a-comprehensive-report}

\begin{center}\rule{0.5\linewidth}{0.5pt}\end{center}

\textbf{Prepared:} July 2026 \textbf{Subject:} Sir Michael Francis
Atiyah (b. April 22, 1929, London, England) \textbf{Position:}
Professor, Oxford University; Director, Isaac Newton Institute
\textbf{Doctoral Advisor:} William V. D. Hodge (Cambridge University)
\textbf{Fields Medal Citation:} \emph{``Did joint work with Hirzebruch
in K-theory; proved jointly with Singer the index theorem of elliptic
operators on complex manifolds; worked in collaboration with Bott to
prove a fixed point theorem related to the `Lefschetz formula'.''}

\begin{center}\rule{0.5\linewidth}{0.5pt}\end{center}

\section{1. Introduction}\label{ch1966michaelfrancisatiyah:introduction}

Michael Atiyah was awarded the Fields Medal in 1966 for his
significant contributions to algebraic topology, K-theory, and the
index theory of elliptic operators. Along with Friedrich Hirzebruch, he
created \textbf{topological K-theory}, a generalized cohomology theory
based on vector bundles that revolutionized algebraic topology. With
Isadore Singer, he proved the \textbf{Atiyah--Singer Index Theorem}, one
of the most significant results of 20th-century mathematics, connecting
analysis, topology, and geometry.

Atiyah's work is characterized by notable breadth and a gift for
forging connections between disparate fields. He made fundamental
contributions to algebraic geometry, gauge theory, mathematical physics,
and the theory of invariants of manifolds. His collaboration with Raoul
Bott produced the \textbf{Atiyah--Bott fixed point theorem}, a powerful
generalization of the Lefschetz fixed point formula. Later in his
career, he turned to questions in theoretical physics, including
instantons, monopoles, and the geometry of gauge theories.

Atiyah's major contributions include: - \textbf{Topological K-theory}
(with Hirzebruch): A generalized cohomology theory using vector bundles.
- \textbf{The Atiyah--Singer Index Theorem}: A formula for the index of
elliptic operators in terms of topological invariants. -
\textbf{Equivariant K-theory} (with Segal): K-theory for group actions.
- \textbf{The Atiyah--Bott fixed point theorem}. - \textbf{The
Atiyah--Hitchin--Singer work on instantons}: Self-dual Yang--Mills
connections. - \textbf{K-theory and the Hopf invariant problem}: A
solution to the classical Hopf invariant one problem. - \textbf{The
Atiyah--Patodi--Singer index theorem}: Index theory for manifolds with
boundary.

This report provides a comprehensive overview of Atiyah's life and his
monumental contributions.

\begin{center}\rule{0.5\linewidth}{0.5pt}\end{center}

\section{2. Biographical Background}\label{ch1966michaelfrancisatiyah:biographical-background}

\subsection{2.1 Early Life and
Education}\label{ch1966michaelfrancisatiyah:early-life-and-education}

Michael Francis Atiyah was born in London on April 22, 1929. His father
was a Lebanese Christian and his mother was Scottish. He grew up in
Khartoum, Sudan, and later in Cairo, Egypt. He attended Victoria College
in Cairo and then Manchester Grammar School in England. He entered
Trinity College, Cambridge in 1947, where he studied under William
Hodge.

\subsection{2.2 Academic Career}\label{ch1966michaelfrancisatiyah:academic-career}

\begin{itemize}
\tightlist
\item
 \textbf{1955:} Ph.D., Cambridge University, under W. V. D. Hodge.
 Thesis: ``Some Applications of Topological Methods to Algebraic
 Geometry.''
\item
 \textbf{1955--1956:} Institute for Advanced Study, Princeton.
\item
 \textbf{1956--1961:} Lecturer, Cambridge University.
\item
 \textbf{1961--1963:} Reader in Mathematics, Oxford University.
\item
 \textbf{1963--1969:} Savilian Professor of Geometry, Oxford.
\item
 \textbf{1966:} Fields Medal at the ICM in Moscow.
\item
 \textbf{1969--1972:} Professor, Institute for Advanced Study,
 Princeton.
\item
 \textbf{1972--1990:} Royal Society Research Professor, Oxford.
\item
 \textbf{1990--1995:} Master, Trinity College, Cambridge.
\item
 \textbf{1995--1997:} Director, Isaac Newton Institute.
\item
 \textbf{2004:} Abel Prize (with Isadore Singer).
\item
 \textbf{2019:} Died at age 89.
\end{itemize}

\subsection{2.3 Doctoral Work}\label{ch1966michaelfrancisatiyah:doctoral-work}

Atiyah's Ph.D.~thesis applied topological methods to algebraic geometry.
He studied the cohomology of sheaves and used spectral sequences to
prove results about algebraic surfaces and ruled surfaces. The work
already displayed his characteristic synthesis of algebraic and
topological methods.

\begin{center}\rule{0.5\linewidth}{0.5pt}\end{center}

\section{3. Topological K-Theory}\label{ch1966michaelfrancisatiyah:topological-k-theory}

\subsection{3.1 Background}\label{ch1966michaelfrancisatiyah:background}

\textbf{K-theory} originated in Grothendieck's work on the Riemann--Roch
theorem in algebraic geometry. Grothendieck defined the group \(K(X)\)
of a scheme \(X\) as the Grothendieck group of coherent sheaves, or of
vector bundles. Atiyah and Hirzebruch adapted this to topology.

\subsection{3.2 Definition}\label{ch1966michaelfrancisatiyah:definition}

For a compact Hausdorff space \(X\), let \(\text{Vect}(X)\) be the set
of complex vector bundles. The isomorphism classes form a monoid under
direct sum. The \textbf{K-group} \(K(X)\) is the Grothendieck group of
this monoid: formal differences \([E] - [F]\) of vector bundles.

The reduced group \(\widetilde{K}(X)\) is the kernel of
\(K(X) \to K(\text{pt}) = \mathbb{Z}\). \textbf{Bott periodicity} gives
the fundamental relation:

\[\widetilde{K}(X) \cong \widetilde{K}(S^2 X)\]

so \(K(S^n)\) is \(\mathbb{Z}\) for even \(n\) and \(0\) for odd \(n\).

\subsection{3.3 The K-Theory Spectrum}\label{ch1966michaelfrancisatiyah:the-k-theory-spectrum}

Atiyah and Hirzebruch showed that \(K\)-theory satisfies the
Eilenberg--Steenrod axioms for a generalized cohomology theory. The
\(K\)-groups \(K^n(X)\) for \(n \in \mathbb{Z}\) are defined by:

\[K^n(X) = \widetilde{K}(S^n X_+)\]

Bott periodicity implies \(K^{n+2}(X) \cong K^n(X)\), so \(K\)-theory is
a \(\mathbb{Z}/2\)-graded theory.

\subsection{3.4 The Hopf Invariant
Problem}\label{ch1966michaelfrancisatiyah:the-hopf-invariant-problem}

One of the first triumphs of K-theory was the solution of the
\textbf{Hopf invariant one problem}. Adams had proved that the only
spheres \(S^n\) that admit a map of Hopf invariant one are
\(S^1, S^3, S^7\). Atiyah and Adams gave a simpler proof using K-theory
and the Adams operations \(\psi^k\).

The proof uses the exterior power operations
\(\lambda^k: K(X) \to K(X)\) and the Adams operations \(\psi^k\), which
are ring homomorphisms satisfying \(\psi^k(x) \equiv x^k \mod p\). For a
Hopf map \(S^{2n-1} \to S^n\), one derives a contradiction unless
\(n=1,2,4\).

\subsection{3.5 Impact}\label{ch1966michaelfrancisatiyah:impact}

K-theory became an essential tool in algebraic topology. It is the
natural home for the index theorem, provides a framework for the theory
of characteristic classes, and connects to geometric topology through
the study of vector fields on spheres, the \(J\)-homomorphism, and
surgery theory.

\begin{center}\rule{0.5\linewidth}{0.5pt}\end{center}

\section{4. The Atiyah--Singer Index
Theorem}\label{ch1966michaelfrancisatiyah:the-atiyahsinger-index-theorem}

\subsection{4.1 Background}\label{ch1966michaelfrancisatiyah:background-1}

For a linear elliptic differential operator
\(D: C^\infty(E) \to C^\infty(F)\) on a compact manifold \(M\), the
kernel and cokernel are finite-dimensional. The \textbf{index} is:

\[\text{ind}(D) = \dim \ker D - \dim \text{coker} D\]

The Russian analysts (Gelfand, Seeley, et al.) had shown the index is a
homotopy invariant, but no formula was known.

\subsection{4.2 The Theorem}\label{ch1966michaelfrancisatiyah:the-theorem}

The Atiyah--Singer Index Theorem states:

\[\text{ind}(D) = \int_M \text{ch}(\sigma(D)) \, \widehat{A}(TM)\]

Here \(\sigma(D)\) is the symbol of \(D\) (an element of K-theory),
\(\text{ch}\) is the Chern character, and \(\widehat{A}(TM)\) is the
\(\widehat{A}\)-genus (a polynomial in the Pontryagin classes of \(M\)).

In K-theoretic form, the theorem is:

\[\text{ind}(D) = p_!([\sigma(D)])\]

where \(p_!\) is the direct image map in K-theory (the \textbf{Gysin
map} or \textbf{pushforward}).

\subsection{4.3 Examples}\label{ch1966michaelfrancisatiyah:examples}

The index theorem includes as special cases: - The \textbf{Gauss--Bonnet
theorem}: \(\chi(M) = \int_M e(TM)\). - The \textbf{Hirzebruch signature
theorem}: \(\tau(M) = \int_M L(TM)\). - The \textbf{Riemann--Roch
theorem}: \(\chi(M, E) = \int_M \text{ch}(E) \, \text{td}(TM)\). - The
\textbf{Atiyah--Singer formula for the Dirac operator}:
\(\text{ind}(\partial\!\!\!/) = \widehat{A}(M)\).

\subsection{4.4 The Proof}\label{ch1966michaelfrancisatiyah:the-proof}

The proof proceeds in several stages: 1. Show that every elliptic
operator is K-theoretically equivalent to a twisted Dirac operator. 2.
Prove the index theorem for twisted Dirac operators using the heat
kernel method. 3. Show that the analytic index and the topological index
agree for all operators, using excision, multiplicativity, and the Bott
periodicity theorem.

\subsection{4.5 Impact}\label{ch1966michaelfrancisatiyah:impact-1}

The index theorem is one of the important results of 20th-century
mathematics. It unifies analysis, topology, and geometry, and has
applications to: - \textbf{Gauge theory}: Instantons, the
Seiberg--Witten equations. - \textbf{Number theory}: The Kronecker limit
formula. - \textbf{Mathematical physics}: Anomalies in quantum field
theory. - \textbf{Differential geometry}: The Hitchin--Thorpe
inequality.

\begin{center}\rule{0.5\linewidth}{0.5pt}\end{center}

\section{5. Equivariant K-Theory}\label{ch1966michaelfrancisatiyah:equivariant-k-theory}

\subsection{5.1 Definition}\label{ch1966michaelfrancisatiyah:definition-1}

For a compact group \(G\) acting on \(X\), \textbf{equivariant K-theory}
\(K_G(X)\) is defined using \(G\)-vector bundles --- vector bundles
\(E \to X\) with a \(G\)-action such that the projection is
\(G\)-equivariant. The equivariant Bott periodicity theorem states:

\[K_G(X) \cong K_G(S^{2V} X)\]

where \(V\) is a \(G\)-representation and \(S^{2V}\) is the one-point
compactification of \(V \oplus V\).

\subsection{5.2 Atiyah and Segal}\label{ch1966michaelfrancisatiyah:atiyah-and-segal}

Atiyah and Graeme Segal developed equivariant K-theory as a fundamental
tool for studying group actions. The \textbf{Atiyah--Segal completion
theorem} relates \(K_G(EG)\) to the completion of the representation
ring:

\[K^*(BG) \cong \widehat{R(G)}\]

This connects the topology of classifying spaces to representation
theory.

\subsection{5.3 The Localization
Theorem}\label{ch1966michaelfrancisatiyah:the-localization-theorem}

The \textbf{Atiyah--Segal localization theorem} states: for a torus
\(T\) acting on \(X\), localization at the maximal ideal of \(R(T)\)
gives an isomorphism:

\[K_T(X) \otimes_{R(T)} \mathbb{C} \cong
K_T(X^T) \otimes_{R(T)} \mathbb{C}\]

where \(X^T\) is the fixed point set. This is an effective tool for
computing equivariant K-theory in terms of fixed point data.

\begin{center}\rule{0.5\linewidth}{0.5pt}\end{center}

\section{6. The Atiyah--Bott Fixed Point
Theorem}\label{ch1966michaelfrancisatiyah:the-atiyahbott-fixed-point-theorem}

\subsection{6.1 Statement}\label{ch1966michaelfrancisatiyah:statement}

For an elliptic operator \(D\) on \(M\) that is \(G\)-equivariant for a
compact group \(G\), the equivariant index \(\text{ind}_G(D)\) is a
character of \(G\). The Atiyah--Bott fixed point theorem computes this
character in terms of data at the fixed points of \(g \in G\).

If \(g\) has isolated fixed points, the formula is:

\[\text{Tr}(g|_{\ker D}) - \text{Tr}(g|_{\text{coker} D}) =
\sum_{p \in M^g} \frac{\text{Tr}(g|_{\sigma_p})}{|\det(1 - g_*)|}\]

where \(\sigma_p\) is the symbol at \(p\) and \(g_*\) is the action on
the tangent space.

\subsection{6.2 Applications}\label{ch1966michaelfrancisatiyah:applications}

\begin{itemize}
\tightlist
\item
 The \textbf{Lefschetz fixed point formula}: the trace of a map
 \(f: M \to M\) on cohomology equals the sum of local contributions at
 fixed points.
\item
 The \textbf{Woodward--Atiyah--Bott formula}: computes characters of
 representations from fixed point data.
\item
 The \textbf{Weil representation}: the fixed point formula gives a
 geometric construction of the Weil representation of symplectic
 groups.
\end{itemize}

\begin{center}\rule{0.5\linewidth}{0.5pt}\end{center}

\section{7. The Dirac Operator and Spin
Geometry}\label{ch1966michaelfrancisatiyah:the-dirac-operator-and-spin-geometry}

\subsection{7.1 The Dirac Operator}\label{ch1966michaelfrancisatiyah:the-dirac-operator}

The \textbf{Dirac operator} on a spin manifold \(M\) is a first-order
elliptic operator \(D: S \to S\) on the spinor bundle. In local
coordinates, \(D = \sum e_i \cdot \nabla_{e_i}\), where \(\cdot\) is
Clifford multiplication.

The index of the Dirac operator is the \(\widehat{A}\)-genus:

\[\text{ind}(D) = \widehat{A}(M)\]

This is the fundamental example of the index theorem, and it shows that
the \(\widehat{A}\)-genus is an integer for spin manifolds.

\subsection{7.2 The Lichnerowicz
Formula}\label{ch1966michaelfrancisatiyah:the-lichnerowicz-formula}

The \textbf{Lichnerowicz formula} relates the square of the Dirac
operator to the Laplacian:

\[D^2 = \nabla^*\nabla + \frac{1}{4}R\]

where \(R\) is the scalar curvature. If \(R > 0\), then \(D^2 > 0\) and
\(\ker D = 0\), so \(\text{ind}(D) = 0\) and \(\widehat{A}(M) = 0\).
This gives the \textbf{Lichnerowicz obstruction}: a manifold with
positive scalar curvature must have vanishing \(\widehat{A}\)-genus.

\subsection{7.3 The Gromov--Lawson--Atiyah
Theorem}\label{ch1966michaelfrancisatiyah:the-gromovlawsonatiyah-theorem}

Atiyah extended this to prove that if a spin manifold \(M\) has a metric
with positive scalar curvature, then its \(\widehat{A}\)-genus vanishes.
More generally, the \(\alpha\)-invariant (the index of the real Dirac
operator) must vanish. This is the \textbf{Atiyah--Hitchin--Singer
obstruction} to positive scalar curvature.

\begin{center}\rule{0.5\linewidth}{0.5pt}\end{center}

\section{8. The Atiyah--Patodi--Singer Index
Theorem}\label{ch1966michaelfrancisatiyah:the-atiyahpatodisinger-index-theorem}

\subsection{8.1 Manifolds with Boundary}\label{ch1966michaelfrancisatiyah:manifolds-with-boundary}

For a manifold \(M\) with boundary \(Y=\partial M\), the index theorem
requires boundary conditions. Atiyah, Patodi, and Singer introduced
\textbf{global boundary conditions} using the Dirac operator on \(Y\).

\subsection{8.2 The APS Index Theorem}\label{ch1966michaelfrancisatiyah:the-aps-index-theorem}

Let \(D\) be a Dirac operator on \(M\) and let \(B\) be the Dirac
operator on \(Y\). The \textbf{APS index theorem} states:

\[\text{ind}(D) = \int_M \widehat{A}(TM) - \frac{1}{2}(\eta(0) + h)\]

where \(\eta(s)\) is the \textbf{eta invariant} of \(B\) (the
regularized number of positive minus negative eigenvalues) and
\(h = \dim \ker B\).

\subsection{8.3 The Eta Invariant}\label{ch1966michaelfrancisatiyah:the-eta-invariant}

The eta invariant \(\eta(0)\) is a spectral invariant that measures the
asymmetry of the spectrum of \(B\). It has deep connections to: -
\textbf{Geometry}: the Chern--Simons invariants. - \textbf{Number
theory}: the Dedekind eta function. - \textbf{Topology}: the
Atiyah--Singer theorem for signature defects of manifolds with boundary.

\begin{center}\rule{0.5\linewidth}{0.5pt}\end{center}

\section{9. Instantons and Gauge
Theory}\label{ch1966michaelfrancisatiyah:instantons-and-gauge-theory}

\subsection{9.1 Yang--Mills Theory}\label{ch1966michaelfrancisatiyah:yangmills-theory}

The \textbf{Yang--Mills equations} are the Euler--Lagrange equations for
the Yang--Mills functional:

\[\mathcal{YM}(A) = \int_M |F_A|^2 \, d\text{vol}\]

where \(A\) is a connection on a principal bundle and \(F_A\) is its
curvature.

\subsection{9.2 Instantons}\label{ch1966michaelfrancisatiyah:instantons}

\textbf{Instantons} are finite-action solutions of the Yang--Mills
equations on \(\mathbb{R}^4\) (or on a compact 4-manifold). They are
\textbf{self-dual} connections: \(*F_A = F_A\) (or anti-self-dual:
\(*F_A = -F_A\)).

Atiyah, Hitchin, and Singer classified instantons on \(S^4\) using the
Atiyah--Singer index theorem. The moduli space \(\mathcal{M}_k\) of
\(SU(2)\) instantons with instanton number \(k\) has dimension \(8k-3\).

\subsection{9.3 The ADHM Construction}\label{ch1966michaelfrancisatiyah:the-adhm-construction}

With Drinfeld, Hitchin, and Manin, Atiyah developed the \textbf{ADHM
construction}, which gives an explicit description of all instantons on
\(\mathbb{R}^4\) in terms of quaternionic linear algebra. The
construction uses an ADHM data set \((B_1, B_2, I, J)\) satisfying
certain equations.

\subsection{9.4 Impact}\label{ch1966michaelfrancisatiyah:impact-2}

The study of instantons led to: - \textbf{Donaldson theory}: invariants
of 4-manifolds from the moduli spaces of instantons. -
\textbf{Seiberg--Witten theory}: a simpler approach to 4-manifold
invariants using monopoles. - \textbf{String theory}: instantons
describe D-brane dynamics and non-perturbative effects.

\begin{center}\rule{0.5\linewidth}{0.5pt}\end{center}

\section{10. The Atiyah--Hitchin--Singer Work on
Monopoles}\label{ch1966michaelfrancisatiyah:the-atiyahhitchinsinger-work-on-monopoles}

\subsection{10.1 Magnetic Monopoles}\label{ch1966michaelfrancisatiyah:magnetic-monopoles}

A \textbf{magnetic monopole} is a solution of the Yang--Mills--Higgs
equations on \(\mathbb{R}^3\) that is static and finite-energy. Atiyah,
Hitchin, and Singer studied \(SU(2)\) monopoles and proved that the
moduli space of charge \(k\) monopoles is a \(4k\)-dimensional
hyper-Kähler manifold, geodesically complete and non-compact.

\subsection{10.2 The Gibbons--Manton
Metric}\label{ch1966michaelfrancisatiyah:the-gibbonsmanton-metric}

The moduli space of \(k=2\) monopoles is the \textbf{Atiyah--Hitchin
manifold}, a 4-dimensional hyper-Kähler manifold that is asymptotic to
\(\mathbb{R}^4/\mathbb{Z}_2\). It plays a central role in the study of
the dynamics of monopoles.

\begin{center}\rule{0.5\linewidth}{0.5pt}\end{center}

\section{11. The Index Theorem for
Families}\label{ch1966michaelfrancisatiyah:the-index-theorem-for-families}

\subsection{11.1 Statement}\label{ch1966michaelfrancisatiyah:statement-1}

Atiyah and Singer extended the index theorem to \textbf{families of
elliptic operators} parametrized by a compact space \(B\). The index of
a family is an element of \(K(B)\), and the family index theorem states:

\[\text{ind}(D) = p_!([\sigma(D)]) \in K(B)\]

where \(p_!\) is the Gysin map for the fibration.

\subsection{11.2 Applications}\label{ch1966michaelfrancisatiyah:applications-1}

The family index theorem is essential for: - \textbf{The gauge anomaly}
in quantum field theory: the cancellation of anomalies is equivalent to
the vanishing of a certain family index. - \textbf{The Thom isomorphism
in K-theory}: the Gysin map \(p_!\) is given by the index of a family of
Dirac operators. - \textbf{The Bismut--Freed computation} of the anomaly
in string theory.

\begin{center}\rule{0.5\linewidth}{0.5pt}\end{center}

\section{12. Mathematical Physics}\label{ch1966michaelfrancisatiyah:mathematical-physics}

\subsection{12.1 Gauge Theories}\label{ch1966michaelfrancisatiyah:gauge-theories}

Atiyah was a pioneer in applying topology to gauge theories. He
organized the 1983 meeting at the Isaac Newton Institute that brought
together mathematicians and physicists, and he co-authored the
influential book \emph{The Geometry and Dynamics of Magnetic Monopoles}
with Hitchin.

\subsection{12.2 Quantum Field Theory}\label{ch1966michaelfrancisatiyah:quantum-field-theory}

Atiyah studied: - \textbf{Anomalies}: The cancellation of gauge
anomalies in quantum field theory, which he showed is equivalent to the
vanishing of the index of a family of Dirac operators. -
\textbf{Topological quantum field theories} (TQFTs): He axiomatized
TQFTs as functors from the category of cobordisms to vector spaces. This
framework underlies the Jones polynomial, the Chern--Simons theory, and
the Seiberg--Witten invariants.

\subsection{12.3 The Jones Polynomial and Knot
Theory}\label{ch1966michaelfrancisatiyah:the-jones-polynomial-and-knot-theory}

Atiyah's work on the Jones polynomial (which he connected to the
representation theory of quantum groups and to Witten's Chern--Simons
theory) was fundamental in establishing the quantum topology program. He
showed that the Jones polynomial arises as the partition function of
\(SU(2)\) Chern--Simons theory on \(S^3\).

\begin{center}\rule{0.5\linewidth}{0.5pt}\end{center}

\section{13. Bott Periodicity and
K-Theory}\label{ch1966michaelfrancisatiyah:bott-periodicity-and-k-theory}

\subsection{13.1 Bott's Theorem}\label{ch1966michaelfrancisatiyah:botts-theorem}

\textbf{Bott periodicity} states that \(\pi_n(U) \cong \pi_{n+2}(U)\)
for the unitary group, and \(\pi_n(Sp) \cong \pi_{n+4}(Sp)\) for the
symplectic group. Equivalently,
\(\Omega^2 BU \simeq BU \times \mathbb{Z}\).

\subsection{13.2 Atiyah's Proof}\label{ch1966michaelfrancisatiyah:atiyahs-proof}

Atiyah gave a proof of Bott periodicity using the theory of Fredholm
operators. The key idea: the space of Fredholm operators on a Hilbert
space is a classifying space for K-theory. If \(H\) is a Hilbert space,
then \(\text{Fred}(H)\) has the homotopy type of
\(BU \times \mathbb{Z}\). The map
\(\text{Fred}(H) \to \Omega^2\text{Fred}(H)\) is a homotopy equivalence,
which is exactly Bott periodicity.

This proof is elegant and conceptual, and it forms the basis for the
topological approach to K-theory.

\begin{center}\rule{0.5\linewidth}{0.5pt}\end{center}

\section{14. The Atiyah--Singer Theorem for Discrete
Groups}\label{ch1966michaelfrancisatiyah:the-atiyahsinger-theorem-for-discrete-groups}

\subsection{14.1 Statement}\label{ch1966michaelfrancisatiyah:statement-2}

For a discrete group \(\Gamma\) acting freely and properly on a manifold
\(M\), the index theorem can be refined to take values in the reduced
\(C^*\)-algebra \(C^*_r(\Gamma)\).

The \textbf{Atiyah \(L^2\) index theorem} states:

\[\text{ind}_\Gamma(D) = \int_{\Gamma \backslash M}
\widehat{A}(T(\Gamma \backslash M))\]

where the left side is the \textbf{\(L^2\)-index}, a trace-class
operator in the von Neumann algebra of \(\Gamma\).

\subsection{14.2 Applications}\label{ch1966michaelfrancisatiyah:applications-2}

\begin{itemize}
\tightlist
\item
 \textbf{The Novikov conjecture}: the higher signatures of a manifold
 are homotopy invariants. This can be attacked using the
 \(\Gamma\)-index theorem and Kasparov's KK-theory.
\item
 \textbf{The Gromov--Lawson conjecture}: manifolds with positive scalar
 curvature are detected by the \(\Gamma\)-index.
\end{itemize}

\begin{center}\rule{0.5\linewidth}{0.5pt}\end{center}

\section{15. The Baum--Connes
Conjecture}\label{ch1966michaelfrancisatiyah:the-baumconnes-conjecture}

Atiyah's work on the index theorem for discrete groups led directly to
the \textbf{Baum--Connes conjecture}, which states that for a discrete
group \(\Gamma\), the assembly map

\[\mu: K_*^\Gamma(\underline{E}\Gamma) \to K_*(C_r^*(\Gamma))\]

is an isomorphism. Here \(\underline{E}\Gamma\) is the universal proper
\(\Gamma\)-space and \(C_r^*(\Gamma)\) is the reduced group
\(C^*\)-algebra. This conjecture connects topology, geometry, and
operator algebras, and it implies the Novikov conjecture and the
Gromov--Lawson conjecture.

\begin{center}\rule{0.5\linewidth}{0.5pt}\end{center}

\section{16. The Atiyah--Singer Theorem and Number
Theory}\label{ch1966michaelfrancisatiyah:the-atiyahsinger-theorem-and-number-theory}

\subsection{16.1 The Kronecker Limit
Formula}\label{ch1966michaelfrancisatiyah:the-kronecker-limit-formula}

Atiyah used the index theorem to prove the \textbf{Kronecker limit
formula}, which expresses the value of Dedekind zeta functions at
\(s=1\):

\[\zeta_K(s) \sim \frac{2^{r_1}(2\pi)^{r_2}}{|\text{disc}(K)|^{1/2}}
\cdot \frac{1}{s-1} + \text{constant}\]

The constant term involves the \(\eta\)-invariant of the Dirac operator
on the associated symmetric space.

\subsection{16.2 Modular Forms}\label{ch1966michaelfrancisatiyah:modular-forms}

Atiyah studied the relationship between the index theorem and modular
forms. The \(\widehat{A}\)-genus appears in the theory of elliptic
genera, which are modular forms associated to a manifold. The
\textbf{Witten genus} is a formal power series whose coefficients are
integrals of characteristic classes, and it is a modular form for
\(SL(2,\mathbb{Z})\).

\begin{center}\rule{0.5\linewidth}{0.5pt}\end{center}

\section{17. The Atiyah--Singer Theorem in Real
Dimension}\label{ch1966michaelfrancisatiyah:the-atiyahsinger-theorem-in-real-dimension}

\subsection{17.1 Real K-Theory}\label{ch1966michaelfrancisatiyah:real-k-theory}

\textbf{Real K-theory} \(KR(X)\) is a \(\mathbb{Z}/8\)-graded cohomology
theory for spaces with an involution. Atiyah introduced \(KR\)-theory to
study real vector bundles and the index theorem for real operators.

\subsection{17.2 The Real Index Theorem}\label{ch1966michaelfrancisatiyah:the-real-index-theorem}

The index theorem for real elliptic operators takes values in
\(KO\)-theory (the K-theory of real vector bundles). Bott periodicity
for \(KO\) is \(\mathbb{Z}/8\)-graded:

\[\widetilde{KO}(S^n) \cong \mathbb{Z}, \mathbb{Z}/2, \mathbb{Z}/2, 0, \mathbb{Z}, 0, 0, 0\]

for \(n \equiv 0,1,2,3,4,5,6,7 \mod 8\).

This is essential for the study of: - \textbf{Real Dirac operators} and
the \(\alpha\)-invariant. - \textbf{Positive scalar curvature metrics}.
- \textbf{The Atiyah--Hitchin--Singer obstruction}.

\begin{center}\rule{0.5\linewidth}{0.5pt}\end{center}

\section{18. The Atiyah--Segal Completion
Theorem}\label{ch1966michaelfrancisatiyah:the-atiyahsegal-completion-theorem}

\subsection{18.1 Statement}\label{ch1966michaelfrancisatiyah:statement-3}

For a compact Lie group \(G\), the \textbf{Atiyah--Segal completion
theorem} states:

\[K^*(BG) \cong \widehat{R(G)}\]

where \(BG\) is the classifying space of \(G\), \(R(G)\) is the
representation ring, and \(\widehat{R(G)}\) is the completion at the
augmentation ideal.

\subsection{18.2 Applications}\label{ch1966michaelfrancisatiyah:applications-3}

This theorem connects topology and representation theory. It allows the
computation of the K-theory of \(BG\) from representation-theoretic
data, and it is essential for: - The study of vector bundles over
classifying spaces. - The computation of the equivariant K-theory of
\(EG\). - The localization theorem in equivariant K-theory.

\begin{center}\rule{0.5\linewidth}{0.5pt}\end{center}

\section{19. The Atiyah--Bott--Shapiro
Construction}\label{ch1966michaelfrancisatiyah:the-atiyahbottshapiro-construction}

\subsection{19.1 Clifford Modules}\label{ch1966michaelfrancisatiyah:clifford-modules}

Atiyah, Bott, and Shapiro gave a construction of KO-theory using
Clifford algebras. They showed that the KO-theory of a point is given by
the Grothendieck group of Clifford modules:

\[\widetilde{KO}(S^n) \cong
\text{Cliff}_{n}/\text{Cliff}_{n-1}\]

where \(\text{Cliff}_n\) is the Grothendieck group of modules over the
Clifford algebra \(C_n\).

\subsection{19.2 The Periodicity}\label{ch1966michaelfrancisatiyah:the-periodicity}

The Clifford algebra \(C_n\) is isomorphic to the \(2 \times 2\) matrix
algebra over \(\mathbb{R}, \mathbb{C}, \mathbb{H}\), or their direct
sums, in a pattern that repeats with period 8. This gives a purely
algebraic proof of Bott periodicity for KO-theory.

\begin{center}\rule{0.5\linewidth}{0.5pt}\end{center}

\section{20. Atiyah's Later Work}\label{ch1966michaelfrancisatiyah:atiyahs-later-work}

\subsection{20.1 Quaternionic Geometry}\label{ch1966michaelfrancisatiyah:quaternionic-geometry}

Atiyah studied quaternionic geometry and its relationship to
hyper-Kähler manifolds and moment maps. He showed that the moduli spaces
of monopoles and instantons are hyper-Kähler, which is essential for
their metric properties.

\subsection{20.2 The Skyrme Model}\label{ch1966michaelfrancisatiyah:the-skyrme-model}

Atiyah's work on the Skyrme model in nuclear physics described baryons
as topological solitons (skyrmions). He showed that the topological
charge of a skyrmion corresponds to the baryon number, and used the
index theorem to prove that the skyrmion moduli space is a hyper-Kähler
manifold.

\begin{center}\rule{0.5\linewidth}{0.5pt}\end{center}

\section{21. Awards and Recognition}\label{ch1966michaelfrancisatiyah:awards-and-recognition}

\textbf{1966:} - Fields Medal (International Mathematical Union).

\textbf{1968:} - Royal Medal of the Royal Society.

\textbf{1983:} - De Morgan Medal, London Mathematical Society.

\textbf{1988:} - Copley Medal of the Royal Society.

\textbf{2004:} - Abel Prize (with Isadore M. Singer).

\textbf{2005:} - Order of Merit (United Kingdom).

\textbf{Elected to:} - Royal Society (1962). - National Academy of
Sciences (USA), Foreign Associate. - French Academy of Sciences, Foreign
Member. - Swedish Royal Academy of Sciences. - Russian Academy of
Sciences, Foreign Member. - American Academy of Arts and Sciences.

\textbf{Knighthood:} - Knighted in 1983. - Member of the Order of Merit
(2005).

\begin{center}\rule{0.5\linewidth}{0.5pt}\end{center}

\section{22. Selected Major
Publications}\label{ch1966michaelfrancisatiyah:selected-major-publications}

\begin{enumerate}
\def\labelenumi{\arabic{enumi}.}
\item
 \textbf{``K-theory.''} \emph{W. A. Benjamin}, 1967. --- \emph{The
 definitive text on topological K-theory.}
\item
 \textbf{``The index of elliptic operators I--V.''} (with I. M. Singer)
 \emph{Annals of Mathematics} 87, 484--530, 546--604, 1968; 88,
 138--172, 1969; 93, 119--138, 1971, 139--149, 1971. --- \emph{The
 Atiyah--Singer index theorem and its generalizations.}
\item
 \textbf{``A Lefschetz fixed point formula for elliptic complexes
 I--II.''} (with R. Bott) \emph{Bulletin AMS} 72, 245--250, 1966;
 \emph{Topology} 7, 93--113, 1968. --- \emph{The Atiyah--Bott fixed
 point theorem.}
\item
 \textbf{``On the structure of instantons.''} (with N. Hitchin, V.
 Drinfeld, Yu. Manin) \emph{Physics Letters A} 65, 185--187, 1978. ---
 \emph{The ADHM construction of instantons.}
\item
 \textbf{``Self-duality in four-dimensional Riemannian geometry.''}
 (with N. Hitchin, I. M. Singer) \emph{Proceedings of the Royal Society
 A} 362, 425--461, 1978. --- \emph{Classification of instantons and the
 index theorem.}
\item
 \textbf{``Spectral asymmetry and Riemannian geometry I--III.''} (with
 V. K. Patodi, I. M. Singer) \emph{Math. Proc. Cambridge Phil. Soc.}
 77, 43--69, 1975; 78, 405--432, 1975; 79, 71--99, 1976. --- \emph{The
 APS index theorem and the eta invariant.}
\item
 \textbf{``Equivariant K-theory and completion.''} (with G. B. Segal)
 \emph{Journal of Differential Geometry} 3, 1--18, 1969. --- \emph{The
 Atiyah--Segal completion theorem.}
\item
 \textbf{``Vector bundles and homogeneous spaces.''} (with F.
 Hirzebruch) \emph{Proceedings of Symposia in Pure Mathematics} 3,
 7--38, 1961. --- \emph{Foundations of topological K-theory.}
\item
 \textbf{``The geometry and dynamics of magnetic monopoles.''} (with N.
 J. Hitchin) \emph{Princeton University Press}, 1988. ---
 \emph{Monopole moduli spaces and hyper-Kähler geometry.}
\item
 \textbf{``The Dirac operator and the Hopf invariant.''} (with J. F.
 Adams) \emph{Topology} 5, 21--71, 1966. --- \emph{K-theoretic solution
 of the Hopf invariant one problem.}
\item
 \textbf{``Riemann surfaces and spin structures.''} \emph{Annales
 Scientifiques de l'ENS} 4, 47--62, 1971. --- \emph{Spin structures on
 Riemann surfaces and theta functions.}
\item
 \textbf{``New invariants of 3- and 4-dimensional manifolds.''} (with
 L. Jeffrey) \emph{The Mathematical Heritage of Hermann Weyl},
 285--299, 1988. --- \emph{Applications of the index theorem to
 low-dimensional topology.}
\end{enumerate}

\begin{center}\rule{0.5\linewidth}{0.5pt}\end{center}

\section{23. References and Further
Reading}\label{ch1966michaelfrancisatiyah:references-and-further-reading}

\subsection{Official Sources}\label{ch1966michaelfrancisatiyah:official-sources}

\begin{itemize}
\tightlist
\item
 IMU Fields Medal 1966:
 https://www.mathunion.org/imu-awards/fields-medal/fields-medals-1966
\item
 Abel Prize 2004: https://www.abelprize.no/abelprize/2004/
\end{itemize}

\subsection{Key Technical References}\label{ch1966michaelfrancisatiyah:key-technical-references}

\begin{itemize}
\tightlist
\item
 Atiyah, M. \emph{Collected Works} (6 vols.). Oxford University Press,
 1988.
\item
 Atiyah, M. \emph{K-Theory}. W. A. Benjamin, 1967.
\item
 Atiyah, M. and Singer, I. M. ``The Index of Elliptic Operators I--V.''
 \emph{Annals of Mathematics}, 1968--1971.
\item
 Atiyah, M. \emph{The Geometry and Physics of Knots}. Cambridge, 1990.
\end{itemize}

\subsection{Textbooks and Surveys}\label{ch1966michaelfrancisatiyah:textbooks-and-surveys}

\begin{itemize}
\tightlist
\item
 Lawson, H. B. and Michelsohn, M.-L. \emph{Spin Geometry}. Princeton
 University Press, 1989.
\item
 Berline, N., Getzler, E., and Vergne, M. \emph{Heat Kernels and Dirac
 Operators}. Springer, 1992.
\item
 Karoubi, M. \emph{K-Theory: An Introduction}. Springer, 1978.
\end{itemize}

\subsection{Biography and Exposition}\label{ch1966michaelfrancisatiyah:biography-and-exposition}

\begin{itemize}
\tightlist
\item
 Cartan, H. ``The Work of Michael Atiyah.'' ICM 1966 Laudatio.
\item
 Hitchin, N. ``Michael Atiyah: A Profile.'' \emph{Abel Prize Laudatio},
 2004.
\item
 Segal, G. ``Michael Atiyah's Work in Algebraic Topology.''
 \emph{Bulletin AMS} 58, 2021.
\item
 Donaldson, S. ``Michael Atiyah and Gauge Theory.'' \emph{Biographical
 Memoirs of the Royal Society}, 2020.
\end{itemize}

\begin{center}\rule{0.5\linewidth}{0.5pt}\end{center}

\emph{(End of report)}
